\scanpage{250}
\mspage{124}
% archive sheet number: 249

\deleted{Corollaire 9.5} \unsure{,}

D'après 9.3, ces conditions sont satisfaites moyennant Loc $(1,2,3,\inserted{4,}5c')$, mais $5c'$ est un marteau-pilon pour une mouche !

On aurait envie \deleted{\unsure{de f}} plutôt d'ajouter \deleted{\unsure{les propriétés}} un axiome aux axiomes Loc $(1)$ à Loc $(\deleted{\unsure{6}}7)$, qui avec Loc $(7)$ serait \unsure{l'axiome} le plus anodin de tous

% slanted
\margin{Il peut être commode \unsure{aussi} d'introduire les conditions encore plus faibles \uline{Loc} $(6\,\unsure{a})$ \deleted{Si $f_1,f_2$ \unsure{sont dans $W$}}, alors $f_1\amalg f_2$ \unsure{est} $\in W$. C'est en fait la condition \unsure{b $1^\circ$} de la prop.~9.4. \uline{Loc} $(6\,b)$ \unsure{Soit|Si} $f_1\amalg f_2$ est dans $W$, alors $f_1$ et $f_2$ \unsure{le sont}. $\Longleftrightarrow$ Soit $Y_0\subset Y$ un summand direct, alors $Y_0\hookrightarrow Y$ est \deleted{\unsure{asph}} dans $\mathrm{Fib}_W$.}

\uline{Loc} $(\deleted{\unsure{6|7}}5)$ \inserted{(Axiome de décomposition, forme faible)} : Si $(\unsure{f^n|f_n|f_i|f^i})_{\unsure{n|i}\in I}$ est une famille finie de flèches de $\mathrm{Cat}$, alors $f=\coprod f_i$ est dans $W$ * \unsure{ssi|ss} les $f_i$ le sont. (Ce sont donc les cond. équiv. de 9.4.)

\uline{Loc} $(\deleted{\unsure{6 bis|7 bis}}5\,\mathrm{bis})$ \inserted{(Axiome de décomp., forme forte)} : comme \uline{Loc} $(\unsure{8|5})$, mais sans supposer $I$ fini.

\uline{Loc} $(5)$ \quad a) Si $f_1,f_2\in W$, alors $f_1\amalg f_2\in W$
\par b) Si $f_1\amalg f_2\in W$ alors $f_1,f_2\in W$
\par ($\Longleftrightarrow$ b$'$) Soit $Y_0\hookrightarrow Y$ summ. direct,
\unsure{alors} $Y_0\hookrightarrow Y$ est dans $\mathrm{Fib}_W$)

\uline{Loc} $(5\,\mathrm{bis})$ comme Loc $(5)$ mais a) renforcé \inserted{(a bis)} pour une famille \inserted{(évent. infinie)} indexée par $I$ arbitraire \unsure{de flèches}.

Ainsi, moyennant Loc $(\deleted{\unsure{6}}7)$, les deux axiomes précédents sont équivalents. De plus, notons ceci :

\deleted{Proposition} \uline{Corollaire} 9.5. (à prop.~9.4.) L'axiome Loc \deleted{$(\unsure{6\,\mathrm{bis}})$}\inserted{$(5\,\mathrm{bis})$} équivaut \deleted{à la}\inserted{à chacune des} propriétés suivantes : \deleted{\unsure{Soit les}}

\textcircled{a} \inserted{Pour tte} $Y$ de $\mathrm{Cat}$, \unsure{décomposé|décomposée} en ses composantes connexes $Y_i$ ($i\in I=\pi_0(Y)$), \deleted{pour} et pour tte flèche $f:X\unsure{-|\to}Y$ dans $\mathrm{Cat}$, on a $f\in W$ \unsure{ss|ssi} les $f_i:X_i=X\mid Y_i-Y_i$ sont $\in W$.

\textcircled{b} Pour $Y,Y_i$ comme dans \textcircled{a}, pour tte
